% !TEX root = main.tex
\documentclass[main.tex]{subfiles}
\begin{document}

\section{Concavity and Null-Spaces of $\jointdual$ and $\centroiddual$}
\label{sec:appendix-nullspace}

The Hessians of $\centroiddual$ and $\jointdual$ derived in Appendices~\ref{sec:appendix-centroid} and~\ref{sec:appendix-joint}, respectively, share the same facet-based structure, from which we derive their concavity and their null-spaces.
We treat $\jointdual$ first and then $\centroiddual$, which corresponds to the special case of fixed weights.
% 
For conciseness, we stack the two multipliers of a site into $\mult_i = [\shift_i, \weight_i] \in \R^{d+1}$.
We then define the affine field
\begin{equation}
  \psi_i(\xx) = 2\,\shift_i\tr (\xx - \tcentroid_i) + \weight_i = \mult_i\tr \phi_i(\xx) ,
  \qquad
  \phi_i(\xx) = \begin{bmatrix} 2(\xx - \tcentroid_i) \\ 1 \end{bmatrix} ,
  \label{eq:appendix-potential}
\end{equation}
through which $\mult_i$ enters the shifted power distance, since, by \eqref{eq:complete-square}, we have $\smash{\norm{\xx \!-\! \qq_i}^2 \!-\! \psi_i(\xx) \!=\! \norm{\xx \!-\! \dual_i}^2 \!-\! \combweight_i}$.
Given that $\psi_i$ is linear in $\mult_i$, a perturbation $\delta\mult_i \!=\! [\delta\shift_i, \delta\weight_i]$ changes it by
\begin{equation}
  \delta\psi_i(\xx) = \delta\mult_i\tr \phi_i(\xx) = 2\,\delta\shift_i\tr (\xx - \tcentroid_i) + \delta\weight_i.
  \label{eq:appendix-dpsi}
\end{equation}
For $\centroiddual$, the weights are fixed, so $\delta\weight_i = 0$ and $\delta\psi_i(\xx) = 2\,\delta\shift_i\tr (\xx - \tcentroid_i)$.

\paragraph{Hessian}
With this notation, the Hessian blocks \eqref{eq:appendix-blocks-off}--\eqref{eq:appendix-blocks-diag} of $\jointdual$ are the entries of the facet-based Laplacian assembled from $\phi_i$:
\begin{subequations}
  \label{eq:appendix-block}
  \begin{empheq}[left={\empheqlbrace\,}]{align}
    \partial^2_{\mult_i \mult_j} \jointdual
      &= \frac{1}{2}\frac{1}{\norm{\dual_i - \dual_j}} \int_{\face_{ij}} \phi_i(\xx)\, \phi_j(\xx)\tr \, dA,
    \label{eq:appendix-block-off}
    \\[2mm]
    \partial^2_{\mult_i \mult_i} \jointdual
      &= -\!\sum_{j \sim i} \frac{1}{2}\frac{1}{\norm{\dual_i - \dual_j}} \int_{\face_{ij}} \phi_i(\xx)\, \phi_i(\xx)\tr \, dA,
    \label{eq:appendix-block-diag}
  \end{empheq}
\end{subequations}
whose shift--shift part is also the Hessian \eqref{eq:centroid-hessian} of $\centroiddual$.

\paragraph{Concavity}
Contracting the Hessian with a perturbation $\delta\mult = \{\delta\mult_i\}$ and pairing the terms of each facet gives
\begin{equation}
  \delta\mult\tr \nabla^2 \jointdual \, \delta\mult
    =
    -\sum_{\{i,j\}}\!\! \frac{1}{2\,\norm{\dual_i \!-\! \dual_j}} \int_{\face_{ij}}\!\! \bigl( \delta\psi_i(\xx) \!-\! \delta\psi_j(\xx) \bigr)^2 dA
    \le 0 ,
  \label{eq:appendix-nsd}
\end{equation}
where the sum runs over pairs of neighbors.
Hence, $\nabla^2 \jointdual \preceq 0$, and restricting \eqref{eq:appendix-nsd} to $\delta\weight = 0$ gives $\nabla^2 \centroiddual \preceq 0$, so both duals are concave.

\paragraph{Homothetic modes of $\jointdual$}
Based on \eqref{eq:appendix-nsd}, the perturbations in \eqref{eq:appendix-dpsi} lie in the null-space of either Hessian if and only if
\begin{equation}
  \delta\psi_i(\xx) = \delta\psi_j(\xx) \quad \forall \xx \in \face_{ij}, \ \forall i \sim j .
  \label{eq:appendix-null-condition}
\end{equation}
% 
To characterize the solutions of \eqref{eq:appendix-null-condition}, we use the lifting of \S\ref{sec:background}, now applied to the shifted weighted sites \eqref{eq:shifted-site}.
Each generator $\dual_i$ is lifted to the point $\smash{(\dual_i, \height_i) \in \R^{d+1}}$ with height $\smash{\height_i \!=\! \norm{\dual_i}^2 \!-\! \combweight_i}$.
The cell $\cell_i$ in \eqref{eq:power-cell} can then be rewritten as every point in $\domain$ that maximizes
\begin{equation}
  a_i(\xx) = 2\,\dual_i\tr\xx - \height_i = 2\,\qq_i\tr\xx - \norm{\qq_i}^2 + \psi_i(\xx).
  \label{eq:appendix-affine}
\end{equation}
Since the sites $\qq_i$ are fixed, the function $a_i(\xx)$ in \eqref{eq:appendix-affine} depends linearly on $\mult_i$ through $\psi_i$, so any perturbation $\delta\mult$ changes $a_i(\xx)$ exactly by $\smash{\delta a_i(\xx) \!=\! \delta\psi_i(\xx)}$.
Geometrically, condition \eqref{eq:appendix-null-condition} requires the perturbations $\delta\psi_i$ and $\delta\psi_j$ of every pair of neighboring cells to agree on their shared facet $\face_{ij}$, and hence on the $(d\!-\!1)$-dimensional hyperplane spanned by it.
% 
Matching the linear and constant parts of $\delta a_i(\xx) = 2\,\delta\dual_i\tr\xx \!-\! \delta\height_i$ with \eqref{eq:appendix-dpsi}, the lifted generators move by
\begin{equation}
  \delta\dual_i = \delta\shift_i ,
  \qquad
  \delta\height_i = 2\,\tcentroid_i\tr\delta\shift_i - \delta\weight_i .
  \label{eq:appendix-lifted-perturbation}
\end{equation}
% 
Now recall from \eqref{eq:lifted-homothety} that a homothety maps every lifted generator to $\smash{s\,(\dual_i, \height_i) + (\translation, c)}$, with the same dilation $s \!>\! 0$ and offset $(\translation, c)$, and thus maps every affine function $a_i(\xx)$ to $\smash{s\,a_i(\xx) + 2\,\translation\tr\xx - c}$ without changing the cells.
An \emph{infinitesimal homothety} is the first-order variation of these maps about the identity $(s, \translation, c) = (1, 0, 0)$.
Writing $s \!=\! 1 + \mu$, each infinitesimal homothety takes the form
\begin{equation*}
  \delta a_i(\xx) = \mu\,a_i(\xx) + 2\,\translation\tr\xx - c ,
\end{equation*}
which satisfies \eqref{eq:appendix-null-condition} since $\delta a_i(\xx) - \delta a_j(\xx) \!=\! \mu\,(a_i(\xx) \!-\! a_j(\xx))$ vanishes on $\face_{ij}$.
Matching the linear and constant parts, the lifted generators move by
\begin{equation}
  \delta\dual_i = \mu\,\dual_i + \translation ,
  \qquad
  \delta\height_i = \mu\,\height_i + c.
  \label{eq:infinitesimal-homothety}
\end{equation}
The corresponding multipliers follow from \eqref{eq:appendix-lifted-perturbation}, which gives $\delta\shift_i = \delta\dual_i$, and from \eqref{eq:appendix-dpsi}, whose shift term vanishes at $\tcentroid_i$, so that $\delta\weight_i = \delta\psi_i(\tcentroid_i) = \delta a_i(\tcentroid_i)$.
Evaluating the infinitesimal homothety at $\tcentroid_i$ with $\smash{a_i(\tcentroid_i) = \gamma_i}$ then yields
\begin{equation}
  \delta\shift_i = \mu\,\dual_i + \translation ,
  \qquad
  \delta\weight_i = \mu\,\gamma_i + 2\,\translation\tr\tcentroid_i - c ,
  \label{eq:appendix-homothety-modes}
\end{equation}
which is the $(d\!+\!2)$-dimensional family \eqref{eq:joint-nullspace} in the null-space of $\nabla^2\jointdual$.
These modes are always present, because the free weights absorb the weight changes required by any homothety.

\paragraph{Homothetic modes of $\centroiddual$}
For $\centroiddual$, the weights are fixed, so a homothety is realizable only if $\delta\weight_i = 0$ for all $i$, \ie only if
\begin{equation}
  \mu\,\gamma_i + 2\,\translation\tr\tcentroid_i - c = 0 \quad \forall i ,
  \label{eq:appendix-centroid-kernel}
\end{equation}
which recovers \eqref{eq:centroid-kernel}.
Since distinct generators make $\delta\shift \!=\! 0$ force $\mu \!=\! 0$ and $\translation = 0$, and then $c \!=\! 0$, distinct triplets $(\mu, \translation, c)$ give distinct modes.
As a result, $\centroiddual$ has a nontrivial homothetic mode if and only if \eqref{eq:appendix-centroid-kernel} has a nontrivial solution, \ie if and only if the vectors $\smash{(\gamma_i, \tcentroid_i, 1)}$ do not span $\R^{d+2}$.
% 
Equivalently, the pivots $\smash{(\tcentroid_i, \gamma_i) \in \R^{d+1}}$ lie on a common hyperplane, as stated in \S\ref{sec:centroid-constraints}.
Since $n$ points in $\R^{d+1}$ always lie on a common hyperplane when $n \!\le\! d\!+\!1$, non-coplanar pivots require $n \!\ge\! d\!+\!2$.

\paragraph{Combinatorial modes}
It remains to discuss whether the homotheties are the only solutions of \eqref{eq:appendix-null-condition}.
Since the supporting hyperplane of facet $\face_{ij}$ is the set of points $\xx$ with $a_i(\xx) \!=\! a_j(\xx)$, an affine function vanishes on $\face_{ij}$ if and only if it is a multiple of $a_i(\xx) \!-\! a_j(\xx)$.
The condition in \eqref{eq:appendix-null-condition} is thus equivalent to
\begin{equation}
  \delta a_j(\xx) - \delta a_i(\xx) = \kappa_{ij}\,(a_j(\xx) - a_i(\xx)) \quad \forall i \sim j ,
  \label{eq:appendix-kappa}
\end{equation}
for some scalars $\kappa_{ij} \!=\! \kappa_{ji}$, one per pair of neighbors.
In terms of the lifted generators $(\dual_i, \height_i)$, \eqref{eq:appendix-kappa} reads
\begin{equation}
  (\delta\dual_j - \delta\dual_i,\, \delta\height_j - \delta\height_i) = \kappa_{ij}\,(\dual_j - \dual_i,\, \height_j - \height_i) ,
\end{equation}
\ie a null-space perturbation may stretch each lifted edge but not rotate it.
These edges connect neighboring generators and form the graph dual to the power diagram, and their projections $\dual_j - \dual_i$ are orthogonal to the facets $\face_{ij}$.
In the following, we characterize the admissible stretch factors $\kappa_{ij}$, which identify the additional null-space modes arising from the combinatorics of the dual graph.

\paragraph{Closed 1-forms}
To this end, we extend the work of \citet{deGoes2013Masonry}, which parametrized the orthogonal duals of a 2D triangulation by closed discrete 1-forms, to power diagrams in $\R^d$.
In our setting, we define $\smash{\omega_{ij} = \kappa_{ij}\,(\dual_j - \dual_i,\, \height_j - \height_i)}$ as a $(d+1)$-sized vector-valued discrete 1-form on the dual graph of our power diagram.
As in \citet{deGoes2013Masonry}, a choice of $\kappa$ is admissible if $\omega$ is \emph{closed}, \ie it sums to zero around the elementary cycles of the dual graph, one around each face of codimension two of the power diagram, where three or more cells meet.
Since we restrict to a simply connected domain $\domain$ (\S\ref{sec:background}) with connected cells $\cell_i$, these elementary cycles generate all cycles of the dual graph, so no harmonic 1-forms exist and every closed $\omega$ is also \emph{exact}, \ie $\smash{\omega_{ij} = (\delta\dual_j - \delta\dual_i,\, \delta\height_j - \delta\height_i)}$.
Integrating $\omega$ thus yields a perturbation of the lifted generators, unique up to a common offset in $\R^{d+1}$.
The null-space of $\nabla^2\jointdual$ is thus parametrized by the scalars $\kappa_{ij}$ for which $\omega$ is closed, together with the $d\!+\!1$ offsets, while the null-space of $\nabla^2\centroiddual$ is its subspace with $\delta\weight_i = \delta a_i(\tcentroid_i) = 0$ for all $i$.

\paragraph{Blocks}
For a power diagram with generators in general position, every face of codimension two of the power diagram (\eg a vertex in 2D, an edge in 3D) is shared by exactly three cells $(i,j,k)$, whose generators span a triangle of the dual triangulation.
Their generators are not collinear, since the facets $\face_{ij}$, $\face_{jk}$ and $\face_{ki}$ would otherwise be parallel and could not meet along a common face, so the lifted edges $\smash{(\dual_j - \dual_i,\, \height_j - \height_i)}$ and $\smash{(\dual_k - \dual_j,\, \height_k - \height_j)}$ are linearly independent.
The closedness condition $\smash{\omega_{ij} + \omega_{jk} + \omega_{ki} = 0}$ then reduces to
\begin{equation}
  (\kappa_{ij} - \kappa_{ki})\,(\dual_j - \dual_i,\, \height_j - \height_i) + (\kappa_{jk} - \kappa_{ki})\,(\dual_k - \dual_j,\, \height_k - \height_j) = 0 ,
\end{equation}
which forces $\kappa_{ij} \!=\! \kappa_{jk} \!=\! \kappa_{ki}$.
Therefore, $\kappa$ is constant over every set of dual edges connected through dual triangles.
Since the dual 
\begin{wrapfigure}[4]{r}{0.45\columnwidth}
  \centering
  \resizebox{\linewidth}{!}{%% generated by figures/nullspace_insets.py; do not edit by hand
\begin{tikzpicture}[x=0.68cm,y=0.68cm]
  \fill[nsBlue!18] (0.000,0.000) -- (1.110,0.000) -- (1.119,0.256) -- (0.573,1.000) -- (0.000,1.000) -- cycle;
  \fill[nsBlue!18] (1.110,2.000) -- (0.000,2.000) -- (0.000,1.000) -- (0.573,1.000) -- (1.119,1.744) -- cycle;
  \fill[nsBlue!18] (1.119,1.744) -- (0.573,1.000) -- (1.119,0.256) -- (1.757,1.000) -- cycle;
  \fill[nsBlue!18] (1.110,0.000) -- (2.660,0.000) -- (2.138,1.000) -- (1.757,1.000) -- (1.119,0.256) -- cycle;
  \fill[nsBlue!18] (2.660,2.000) -- (1.110,2.000) -- (1.119,1.744) -- (1.757,1.000) -- (2.138,1.000) -- cycle;
  \fill[nsCut] (2.660,0.000) -- (3.251,0.000) -- (3.801,1.000) -- (3.251,2.000) -- (2.660,2.000) -- (2.138,1.000) -- cycle;
  \fill[nsGreen!18] (3.251,0.000) -- (4.893,0.000) -- (4.282,1.000) -- (3.801,1.000) -- cycle;
  \fill[nsGreen!18] (4.893,2.000) -- (3.251,2.000) -- (3.801,1.000) -- (4.282,1.000) -- cycle;
  \fill[nsGreen!18] (4.893,0.000) -- (5.600,0.000) -- (5.600,2.000) -- (4.893,2.000) -- (4.282,1.000) -- cycle;
  \draw[nsFacet] (0.573,1.000) -- (0.000,1.000);
  \draw[nsFacet] (1.119,0.256) -- (0.573,1.000);
  \draw[nsFacet] (1.110,0.000) -- (1.119,0.256);
  \draw[nsFacet] (0.573,1.000) -- (1.119,1.744);
  \draw[nsFacet] (1.110,2.000) -- (1.119,1.744);
  \draw[nsFacet] (1.119,0.256) -- (1.757,1.000);
  \draw[nsFacet] (1.119,1.744) -- (1.757,1.000);
  \draw[nsFacet] (2.138,1.000) -- (1.757,1.000);
  \draw[nsFacet] (2.660,0.000) -- (2.138,1.000);
  \draw[nsFacet] (2.660,2.000) -- (2.138,1.000);
  \draw[nsFacet] (3.251,0.000) -- (3.801,1.000);
  \draw[nsFacet] (3.801,1.000) -- (3.251,2.000);
  \draw[nsFacet] (4.282,1.000) -- (3.801,1.000);
  \draw[nsFacet] (4.893,0.000) -- (4.282,1.000);
  \draw[nsFacet] (4.893,2.000) -- (4.282,1.000);
  \draw[nsDomain] (0.000,0.000) rectangle (5.600,2.000);
  \draw[nsDual, nsBlue] (0.400,0.450) -- (0.400,1.550);
  \draw[nsDual, nsBlue] (0.400,0.450) -- (1.150,1.000);
  \draw[nsDual, nsBlue] (0.400,0.450) -- (1.850,0.400);
  \draw[nsDual, nsBlue] (0.400,1.550) -- (1.150,1.000);
  \draw[nsDual, nsBlue] (0.400,1.550) -- (1.850,1.600);
  \draw[nsDual, nsBlue] (1.150,1.000) -- (1.850,0.400);
  \draw[nsDual, nsBlue] (1.150,1.000) -- (1.850,1.600);
  \draw[nsDual, nsBlue] (1.850,0.400) -- (1.850,1.600);
  \draw[nsDual, nsBlue] (1.850,0.400) -- (3.000,1.000);
  \draw[nsDual, nsBlue] (1.850,1.600) -- (3.000,1.000);
  \draw[nsDual, nsGreen] (3.000,1.000) -- (4.000,0.450);
  \draw[nsDual, nsGreen] (3.000,1.000) -- (4.000,1.550);
  \draw[nsDual, nsGreen] (4.000,0.450) -- (4.000,1.550);
  \draw[nsDual, nsGreen] (4.000,0.450) -- (4.900,1.000);
  \draw[nsDual, nsGreen] (4.000,1.550) -- (4.900,1.000);
  \fill[nsSite] (0.400,0.450) circle (0.07);
  \fill[nsSite] (0.400,1.550) circle (0.07);
  \fill[nsSite] (1.150,1.000) circle (0.07);
  \fill[nsSite] (1.850,0.400) circle (0.07);
  \fill[nsSite] (1.850,1.600) circle (0.07);
  \fill[nsSite] (3.000,1.000) circle (0.07);
  \fill[nsSite] (4.000,0.450) circle (0.07);
  \fill[nsSite] (4.000,1.550) circle (0.07);
  \fill[nsSite] (4.900,1.000) circle (0.07);
  \node[nsLabel, below=2pt] at (3.000,1.000) {$i$};
\end{tikzpicture}
}
\end{wrapfigure}
triangles generate all dual cycles, these sets are exactly the \emph{blocks} of the dual graph (see inset), \ie its maximal subgraphs that remain connected after removing any single node, and a constant $\kappa$ per block always defines a closed 1-form.
The null-space of $\nabla^2\jointdual$ thus has dimension $b + d + 1$, where $b$ is the number of blocks.
When the dual graph is a single block, \ie $b\!=\!1$, $\kappa \equiv \mu$ and $\smash{\omega_{ij} = \mu\,(\dual_j - \dual_i,\, \height_j - \height_i)}$, so the only solutions are the $d\!+\!2$ homotheties as in \eqref{eq:infinitesimal-homothety}.
% The simplest instance is the volume Hessian \eqref{eq:newton-step}, for which $\delta\shift = 0$ makes each $\delta a_i = \delta\weight_i$ constant, so \eqref{eq:appendix-kappa} forces $\kappa \equiv 0$ and $\omega = 0$, leaving only the common offset $\delta\weight_i = c$.

\paragraph{Cut cells}
When the dual graph has $b \!\ge\! 2$ blocks, the additional $b \!-\! 1$ modes beyond the homotheties come from the arrangement of the cells, since each block carries its own constant $\kappa$.
Distinct blocks only meet at \emph{cut cells}, \ie cells $\cell_i$ whose removal splits $\domain$ into two or more parts, or equivalently dual nodes shared by several blocks, such as the orange cell $i$ in the inset above.
Setting $\delta a_k(\xx) \!=\! 0$ on all but one side of cell $i$ and $\smash{\delta a_k(\xx) \!=\! \kappa\,(a_k(\xx) - a_i(\xx))}$ for every cell $k$ on the remaining side applies an infinitesimal dilation about cell $i$ to that side only.
Each cut cell shared by $m$ blocks thus contributes $m \!-\! 1$ independent modes, one per side beyond the first, which together account for the $b \!-\! 1$ additional modes.

\paragraph{Bridges}
A block may also reduce to a single dual edge between two cut cells, or between a cut cell and a cell with no other neighbor, in which case the facet itself separates the domain.
We thus define a 
\begin{wrapfigure}[5]{r}{0.45\columnwidth}
  \centering
  \vspace{-2mm}
  \resizebox{\linewidth}{!}{%% generated by figures/nullspace_insets.py; do not edit by hand
\begin{tikzpicture}[x=0.68cm,y=0.68cm]
  \fill[nsShade] (0.000,0.000) -- (1.292,0.000) -- (0.908,1.000) -- (0.000,1.000) -- cycle;
  \fill[nsShade] (1.292,2.000) -- (0.000,2.000) -- (0.000,1.000) -- (0.908,1.000) -- cycle;
  \fill[nsShade] (1.292,0.000) -- (2.496,0.000) -- (2.496,2.000) -- (1.292,2.000) -- (0.908,1.000) -- cycle;
  \draw[nsFacet] (0.908,1.000) -- (0.000,1.000);
  \draw[nsFacet] (1.292,0.000) -- (0.908,1.000);
  \draw[nsFacet] (1.292,2.000) -- (0.908,1.000);
  \draw[nsFacet] (2.496,0.000) -- (2.496,2.000);
  \draw[nsFacet] (3.839,1.000) -- (3.482,2.000);
  \draw[nsFacet] (3.482,0.000) -- (3.839,1.000);
  \draw[nsFacet] (5.000,1.000) -- (3.839,1.000);
  \draw[nsDomain] (0.000,0.000) rectangle (5.000,2.000);
  \draw[nsHighlight] (2.496,0.000) -- (2.496,2.000);
  \node[nsLabel, above] at (2.496,2.000) {$\face_{ij}$};
  \draw[nsDual, nsGray] (0.450,0.500) -- (0.450,1.500);
  \draw[nsDual, nsGray] (0.450,0.500) -- (1.750,1.000);
  \draw[nsDual, nsGray] (0.450,1.500) -- (1.750,1.000);
  \draw[nsDual, nsRed] (1.750,1.000) -- (3.050,1.000);
  \draw[nsDual, nsGray] (3.050,1.000) -- (4.450,1.500);
  \draw[nsDual, nsGray] (3.050,1.000) -- (4.450,0.500);
  \draw[nsDual, nsGray] (4.450,1.500) -- (4.450,0.500);
  \fill[nsSite] (0.450,0.500) circle (0.07);
  \fill[nsSite] (0.450,1.500) circle (0.07);
  \fill[nsSite] (1.750,1.000) circle (0.07);
  \fill[nsSite] (3.050,1.000) circle (0.07);
  \fill[nsSite] (4.450,1.500) circle (0.07);
  \fill[nsSite] (4.450,0.500) circle (0.07);
  \node[nsLabel, below left=1pt] at (1.750,1.000) {$i$};
  \node[nsLabel, below right=1pt] at (3.050,1.000) {$j$};
\end{tikzpicture}
}
\end{wrapfigure}
\emph{bridge} as a facet $\face_{ij}$ that alone cuts $\domain$ into two parts, \ie a dual edge that lies on no dual cycle and, therefore, forms a block on its own (see inset).
Its scalar $\kappa_{ij}$ is then free, which is one of the $b \!-\! 1$ modes counted above, and setting $\delta a_k(\xx) = 0$ on the side of cell $j$ and $\smash{\delta a_k(\xx) = \kappa_{ij}\,(a_i(\xx) - a_j(\xx))}$ for every cell $k$ on the side of cell $i$ rotates all the lifted hyperplanes on that side about the facet.
For instance, a cell $i$ with a single neighbor $j$ is separated by the bridge $\face_{ij}$, and $j$ is a cut cell unless it also has no other neighbor.

\paragraph{Strips}
\begin{wrapfigure}[4]{r}{0.42\columnwidth}
\vspace{1mm}
  \centering
  \resizebox{\linewidth}{!}{%% generated by figures/nullspace_insets.py; do not edit by hand
\begin{tikzpicture}[x=0.68cm,y=0.68cm]
  \fill[nsCut] (0.692,0.000) -- (1.556,0.000) -- (1.556,1.200) -- (0.692,1.200) -- cycle;
  \fill[nsCut] (1.556,0.000) -- (2.284,0.000) -- (2.284,1.200) -- (1.556,1.200) -- cycle;
  \fill[nsCut] (2.284,0.000) -- (3.145,0.000) -- (3.145,1.200) -- (2.284,1.200) -- cycle;
  \draw[nsFacet] (0.692,0.000) -- (0.692,1.200);
  \draw[nsFacet] (1.556,0.000) -- (1.556,1.200);
  \draw[nsFacet] (2.284,0.000) -- (2.284,1.200);
  \draw[nsFacet] (3.145,0.000) -- (3.145,1.200);
  \draw[nsDomain] (0.000,0.000) rectangle (3.900,1.200);
  \draw[nsDual, nsBlue] (0.350,0.600) -- (1.100,0.600);
  \draw[nsDual, nsOrange] (1.100,0.600) -- (1.900,0.600);
  \draw[nsDual, nsGreen] (1.900,0.600) -- (2.750,0.600);
  \draw[nsDual, nsPurple] (2.750,0.600) -- (3.500,0.600);
  \fill[nsSite] (0.350,0.600) circle (0.07);
  \fill[nsSite] (1.100,0.600) circle (0.07);
  \fill[nsSite] (1.900,0.600) circle (0.07);
  \fill[nsSite] (2.750,0.600) circle (0.07);
  \fill[nsSite] (3.500,0.600) circle (0.07);
\end{tikzpicture}
}
\end{wrapfigure}
Finally, a \emph{strip} of $n$ parallel cells has a path as its dual graph, whose $n \!-\! 1$ edges are all bridges and whose $n \!-\! 2$ interior cells are all cut cells, so it has $b = n \!-\! 1$ blocks and $n \!-\! 2$ additional modes (see inset).

% For $\nabla^2\centroiddual$, an arrangement mode must also satisfy the pivot conditions $\delta a_i(\tcentroid_i) = 0$; for instance, the rotation of a boundary cell with a single neighbor $j$ about its facet is realizable only if $\tcentroid_i$ lies on the supporting hyperplane $\{\xx : a_i(\xx) = a_j(\xx)\}$.

\end{document}
